% templates/C_flow_ascii_ru.tex
% Русский ASCII-поток для отчётов C (pdfLaTeX-safe)
% Без Unicode-символов Σ/χ/→/⊤/∥: используем "Sigma/chi", "->/<-", "Pi_top", "V_par".
% Требует в преамбуле: \usepackage{fancyvrb}

\begin{figure}[p]
\centering
\thispagestyle{empty}

\fbox{%
\begin{minipage}{0.97\textwidth}
\small
\begin{Verbatim}[fontsize=\footnotesize]
+--------------------------- GEOMETRIYA EVO ----------------------------+
|  Sigma: granichnaya poverhnost (parametrizaciya X(u,v))               |
|  n(x): normal' na Sigma                                               |
|  e_par: prodol'naya os' (kanonicheski iz geometrii; CANON §17)        |
|  chi in {-1,+1}: pasport hiral'nosti (diskretnyy znak)                |
|  s_sh in {-1,0,+1}: rezhim shneka (-> / stoyachiy / <-)               |
+--------------------------------+--------------------------------------+
                                 |
                                 v
+---------------- DVUTORIY: GEOMETRIZACIYA (B->C, KANDIDAT) ------------+
|  q in R\{0}: parametr namotki (q=0 -> tonkiy tor; q!=0 -> dvutoriy)   |
|  w(x): geometricheskiy ves ("vnesh./vnutr. rychag") na Sigma          |
|  w_mode in {binary, hybrid, ...},  w_asym: metrika nesimmetrii w(x)   |
|  Norma karkasa: V_space=0 bez (q!=0 i w_asym>0);                      |
|                pri s_sh=0 (net tor-mody) dreyf zapreshchen.           |
+--------------------------------+--------------------------------------+
                                 |
                                 v
+------------------- PROKSI NA Sigma (CANON §16) -----------------------+
|  K(x)=n <v^2>        G(x)=n <v>        P_ij=n <v_i v_j>              |
|  A_ij=P_ij-(1/3)tr(P) delta_ij     Pi_top(G)=G-(G·n)n               |
|  Proksi — IZMERITEL'NYE YARLYKI (ne sushchnosti, ne "polya").         |
|  Prichinnost': statistika kontaktov / Delta p na Sigma + rezhimy Sigma|
|  (bez "potoka SE v ob'eme").                                          |
+--------------------------------+--------------------------------------+
                                 |
                 +---------------+-----------------+
                 |                                 |
                 v                                 v
+-------------------- j_ext (vneshniy narushitel') ---------------------+
|  j_ext(x)=F_ext(K,G,A; Sigma)                                         |
|  (grad K / prihodyashchaya VR kak perekos statistiki na Sigma;        |
|   bez "potoka SE v ob'eme")                                           |
+--------------------------------+--------------------------------------+

+-------------------- j_int (vnutrenniy: tor -> sp) --------------------+
|  g_top := Pi_top(G)                                                   |
|  a_top := Pi_top(A n)    (gde (A n)_i = A_ij n_j)                     |
|                                                                        |
|  j_tor(x)= chi*s_sh( a1 n x g_top + a2 n x a_top )                     |
|  j_sp(x) = beta(K) w(x) (j_tor · t_theta) h(x)                         |
|  j_int(x)= j_tor(x) + j_sp(x)                                          |
|                                                                        |
|  Garantii: s_sh=0 => j_tor=0 => j_sp=0                                 |
|           q=0 i/ili w~const => integral'naya kompensaciya => V_space~0 |
|  B->C logi (obyazatel'no): q,w_mode,w_asym,j_tor_norm,j_sp_norm        |
+--------------------------------+--------------------------------------+
                                 |
                                 +--------------------------+
                                                            |
                                                            v
+--------------------------- SUMMA NA Sigma ----------------------------+
|  j(x) = j_ext(x) + j_int(x)                                           |
+--------------------------------+--------------------------------------+
                                 |
                                 v
+------------------------ INTEGRAL'NYJ OTKLIK --------------------------+
|  V     = mu_T * int_Sigma j dS                                        |
|  V_par = V · e_par                                                     |
|  Omega = muR_tilde * int_Sigma (r-r0) x j dS                           |
|  J_int     = int_Sigma j_int dS                                       |
|  J_int_par = J_int · e_par                                            |
+--------------------------------+--------------------------------------+
                                 |
                                 v
+---------------------- SHLYUZY / VERIFIKACIYA -------------------------+
|  NULEVOY KORIDOR (CANON §17.2):                                       |
|    v_rms = sqrt(<K/n>_Sigma),  n_bar = <n>_Sigma                      |
|    eps_G   = ||Pi_top(G)||_Sigma / (n_bar v_rms)                      |
|    eps_A   = ||A||_Sigma / (n_bar v_rms^2)                            |
|    eps_V   = ||V||/v_rms,   eps_Vpar = |V_par|/v_rms                  |
|                                                                        |
|  Razvedenie norm:                                                      |
|    - T2: "V=0" — vektornyy nol' -> shlyuz po eps_V                     |
|    - T5-T6 + stoyachiy: norma po V_par -> shlyuz po eps_Vpar           |
|    (v oboih sluchayah "net narushitelya" tol'ko pri eps_G, eps_A v nule)|
|                                                                        |
|  T2 auto-shlyuz: (eps_G, eps_A ~ 0) => (eps_V ~ 0)                    |
|                                                                        |
|  T5-T6 (shnek, §17):                                                   |
|    (1) chi-nechyotnost': V_par(chi)= -V_par(-chi) (parnyy progon)      |
|    (2) revers bega:      V_par(->)= -V_par(<-) (parnyy, fiks. chi)     |
|    (3) stoyachiy: s_sh=0 i "net narushitelya" => V_par ~ 0             |
|                                                                        |
|  OKHRANNIKI GEOMETRIZACII (obyazatel'ny pered D1):                     |
|    N1: q=0, s_sh!=0 -> Omega!=0 dopustimo, no V~0                      |
|    N2: q!=0, s_sh=0  -> V=0                                            |
|    N3: q!=0, w_asym~0 -> V~0                                           |
|  D1: q!=0, w_asym>0, s_sh!=0 -> V_par!=0 + normy §17                  |
+------------------------------------------------------------------------+
\end{Verbatim}
\end{minipage}%
}

\caption{ASCII-potok (RU): Sigma -> (q,w) -> proksi -> j\_ext / j\_int (tor->sp) -> integraly -> shlyuzy (T2,T5--T6) -> okhranniki N1--N3 -> D1.}
\label{fig:C_flow_ascii_ru}
\end{figure}
